% U-Speak Lab ／ 最終要件定義：多読の 島（よみもの島）と こども向け UI（Raz-Kids の 置き換え提案に 向けて）
%   組む:  cd docs && xelatex reading-island-spec.tex
\documentclass[a4paper,10pt]{article}
\usepackage{fontspec}
\usepackage{xeCJK}
\usepackage{tcolorbox}
\usepackage{booktabs}
\usepackage{array}
\usepackage{longtable}
\usepackage{enumitem}
\usepackage{titlesec}
\usepackage{fancyhdr}
\usepackage[margin=18mm,top=22mm,bottom=20mm]{geometry}
\tcbuselibrary{skins,breakable}

\newcommand{\uspeakfont}[1]{\setCJKmainfont{#1}\setmainfont{#1}}
\IfFontExistsTF{Noto Sans CJK JP}{\uspeakfont{Noto Sans CJK JP}}{%
\IfFontExistsTF{Noto Sans JP}{\uspeakfont{Noto Sans JP}}{%
\IfFontExistsTF{Yu Gothic}{\uspeakfont{Yu Gothic}}{%
\IfFontExistsTF{Hiragino Sans}{\uspeakfont{Hiragino Sans}}{%
  \GenericError{}{U-Speak: no Japanese font found}{}{Install Noto Sans CJK JP.}%
}}}}

\definecolor{deep}{HTML}{1B3A2F}
\definecolor{orange}{HTML}{E27A2D}
\definecolor{cream}{HTML}{FBF6EC}
\definecolor{peach}{HTML}{FDEBDD}
\definecolor{ink}{HTML}{22372F}
\definecolor{soft}{HTML}{6B7F74}
\definecolor{rule}{HTML}{CFC9B4}

\setlength{\parindent}{0pt}
\setlength{\parskip}{4pt}
\renewcommand{\arraystretch}{1.35}
\color{ink}
\setlist{nosep,leftmargin=1.4em,itemsep=2pt}

\pagestyle{fancy}
\fancyhf{}
\renewcommand{\headrule}{\color{rule}\hrule height 0.4pt}
\fancyhead[L]{\footnotesize\color{deep}最終要件定義：多読の 島（よみもの島）と こども向け UI}
\fancyhead[R]{\footnotesize\color{deep}合同会社 U-Speak Lab（社内）}
\fancyfoot[C]{\footnotesize\color{soft}\thepage}

\titleformat{\section}{\large\bfseries\color{deep}}{\thesection}{0.6em}{}[{\color{orange}\titlerule[1pt]}]
\titleformat{\subsection}{\normalsize\bfseries\color{deep}}{\thesubsection}{0.5em}{}
\titlespacing*{\section}{0pt}{12pt}{6pt}
\titlespacing*{\subsection}{0pt}{8pt}{3pt}

\newtcolorbox{summary}[1][]{enhanced,breakable,colback=cream,colframe=deep,boxrule=0pt,leftrule=4pt,arc=2pt,left=8pt,right=8pt,top=6pt,bottom=6pt,#1}
\newtcolorbox{note}[1][]{enhanced,breakable,colback=peach,colframe=orange,boxrule=0.6pt,arc=3pt,left=8pt,right=8pt,top=5pt,bottom=5pt,#1}
\newcommand{\req}[1]{\textbf{\color{orange}#1}}
\newcolumntype{L}[1]{>{\raggedright\arraybackslash}p{#1}}

\begin{document}

\begin{tcolorbox}[enhanced,colback=deep,colframe=deep,arc=4pt,left=12pt,right=12pt,top=10pt,bottom=10pt]
{\color{white}\footnotesize FINAL REQUIREMENTS · v1.0（承認待ち）}\\[2pt]
{\color{white}\LARGE\bfseries 多読の 島「よみもの島」と こども向け UI}\\[3pt]
{\color{white}\normalsize こう 開発します ― Raz-Kids の 置き換え提案に 向けた 最終要件定義}
\end{tcolorbox}

{\footnotesize\color{soft}作成：2026 年 11 月 17 日／作成者：U-Speak Lab（Claude）／もとの 材料：「Raz-Kids 徹底リサーチ」（2026-11-17）、商談記録（スマイルキッズ様・まなびば Jr. 様）、U-Speak Web の 現在の 実装\\
\textbf{この 文書が 承認（Yes）されるまで、実装は はじめません。}}

\begin{summary}
\textbf{結論}\quad Raz-Kids の 強みは「\textbf{レベル別の 英語絵本を、聞いて・読んで・クイズに 答えて、ごほうびを 集める}」に しぼった \textbf{迷わない 作り}。
U-Speak は 話す 練習・みんなで 遊べる 島・LINE の 保護者レポートで 勝っている が、\textbf{画面の わかりやすさ}で 負けている。\\[3pt]
そこで ① \textbf{多読の 島「よみもの島」}（オリジナル絵本を ロードマップの 段階ごとに）を 新しく つくり、② \textbf{こども向けの 迷わない UI}（きょうの がくしゅう・全テキスト読み上げ・1 画面 1 問・進み具合バー・入力なしの クラス参加・タップで 移動・保護者の 数字 5 つ・3 ステップの 課題配信）を 入れ、③ それらを まとめて \textbf{先生が 管理画面で オンに できる「幼児モード」}（Raz-Kids なみに かんたんな 画面）に する。
\textbf{U-Speak に もう ある 機能は ぜんぶ 使い}、新しく つくるのは 絵本の 画面・絵本の 材料づくり・こども用ホームの 3 つに しぼる。
\end{summary}

\section{目的と 背景}
\begin{itemize}
\item 商談 6 件のうち 2 件（スマイルキッズ様・まなびば Jr. 様）が Raz-Kids を 導入ずみで、比較を 求められた。どちらも「園が まとめて 契約し 子どもに 配る」予算が ある。スマイルキッズ様は「絵本の 部屋が あれば 置き換えも」。
\item 価格は Raz-Kids（日本の 取り扱い 1 人 年 2,200〜3,300 円）より U-Speak（月 800〜1,200 円）が 数倍 高い。置き換えには \textbf{多読で 同じ ことが できる}＋\textbf{話す・LINE レポート・教室専用}の 上乗せを 見せる 必要が ある。
\item 目標：\textbf{体験の 初日で 離脱しない}・\textbf{迷わない}・\textbf{伸びが 見える}・\textbf{Raz-Kids の 置き換え提案が できる}。
\end{itemize}

\section{開発の 基本方針（ぜんぶの 要件に かかる きまり）}
\begin{enumerate}
\item \req{既存機能を フルに 使う}：レベル＝ロードマップ、ごほうび＝U-Speak コイン と きせかえ・じぶんの 島、クイズ＝問題DB と 7 形式の 部品、読み上げ＝自然な 声（/api/tts）、音読判定＝面接の 音声認識、課題＝宿題、先生＝/class、保護者＝web レポートと LINE。\textbf{新しい レベル・通貨・先生画面は つくらない。}
\item \req{判定は サーバー}：クイズの 正解・コイン・記録・段階アップは サーバーが きめる。答えは client/dist に 置かない。
\item \req{知的財産}：Raz-Kids の 本文・絵・キャラクター・画面デザイン・名前（Star Zone・Raz Rocket など）は 使わない・まねない。「レベル別多読」「1 画面 1 問」「進み具合バー」などの \textbf{しくみ}だけを、自社の 絵・文・デザインで つくる。
\item \req{こどもの 文字}：説明は ひらがな・分かち書き。絵本の 本文・ことば・クイズは 学習の 中身なので 訳さない（\texttt{translate="no"}）。字は 11px 以上・ボタンは 34px 以上。iPad と スマホの 横もちで 1 画面に おさまる。
\item \req{個人情報}：子どもの 音読の 録音は 個人情報。のこす 期間・聞ける 人・1 人ぶんの 削除を きめて 守る（§6）。
\end{enumerate}

\section{機能要件}
\subsection{F1　多読の 島「よみもの島」（本命）}
\begin{longtable}{@{}L{30mm}L{140mm}@{}}
\toprule
島 & 3D の 島（\texttt{island-kit}）。としょかん（まんなか）・きく ベンチ・おんどく ブース・クイズの いえ・よんだ 本の ひろば。しまの ちずの「おすすめ」に 入れる。右の バーの 1 ボタンからも 本だなを ひらける（島に 行かなくても 読める）。\\
本だな & \textbf{ロードマップの 12 段階（A・B・C・1〜9）で 本を ならべる}。その 子の いまの 段階の 本＋1 つ 下の 段階（かんたん）。先生が /class で 段階を かえると 本だなも かわる。表紙・タイトル（英語）・✓（読んだ）・★の 数。\\
1 冊の 流れ & \textbf{① きく}：1 ページずつ 自然な 声で 読み上げ、読んでいる ことばを 光らせる。ゆっくり（0.75 倍）。\newline
\textbf{② よむ}：じぶんで めくる。ことばを おすと その ことばの 発音と 日本語の いみ。\newline
\textbf{③ おんどく}：文を 読むと 音声認識で 読めた ことばが 緑に。正確さ（％）と 1 分の 語数を 出す。\newline
\textbf{④ クイズ}：5 問・1 画面 1 問（F5）。答えは サーバー。\\
✓ と ★ & 4 つの うち 「きく・よむ・クイズ 80\% 以上」で ✓。★ は 4 つの 活動ごとに 1（最大 4）。★ は 本の しるし（通貨では ない）。\\
ごほうび & 本 1 冊 ✓ で U-Speak コイン 20 まい（1 日の 上限 あり）。コインで \textbf{きせかえ・じぶんの 島・ぶきや}（すでに ある）を 買う＝「学習と ごほうびを わける」Raz-Kids の 形を 既存の 店で 実現。\\
段階アップ & \textbf{ロードマップが きめる}。クイズの 答えは 学習記録（\texttt{roblox\_events} の quiz・world=reading・その 段階の 級）に 入り、ロードマップの「その 級の クイズ 20 問・8 割」・英検の 合格見込み・苦手単語の 復習に そのまま つながる。島の めやす「この 段階の 本 ○ 冊」「音読の 正確さ 95\%」は ロードマップ画面に「よんだ 本」として 出す（段階アップの 条件に 足すかは 運用で きめる）。\\
\bottomrule
\end{longtable}

\subsection{F2　絵本の 材料づくり（オリジナル）}
\begin{itemize}
\item \textbf{材料}：AI で つくる オリジナルの 物語（段階ごとの 語数・文の 長さ・語い）、提携している 外国語大学の 素材、CC BY の オープン絵本（本ごとに ライセンス確認・NC は 使わない）。\textbf{市販の 絵本・教材は 使わない}。
\item \textbf{流れ}：① 物語の 下書き（段階の 単語リストに そって）→ ② 絵（既存の 絵づくりの 流れ。登場人物は 本ごとに 同じ 見た目）→ ③ 読み上げ音声（/api/tts で 先に 作って 保存）→ ④ クイズ 5 問の 下書き → ⑤ 運営が たしかめて 公開。
\item \textbf{量}：最初 \textbf{36 冊}（12 段階 × 3 冊）、つぎに \textbf{120 冊}（× 10）。1 冊 4〜16 ページ（段階で かわる）。
\item \textbf{入れ口}：クイズは \textbf{問題DB（/admin/qbank）}に 形式「えほん」で 入れる。本の 一覧・段階・出典・公開の オンオフだけ 管理画面に 1 ページ。PDF・画像・CSV からの 読みとりは 宿題の 読みとり（AI）を 使う。
\item \textbf{データ}：\texttt{\{ id, title, stage, words, source, pages:[\{text, image, audio\}], glossary:[\{en, ja\}] \}}。本文と 絵は \texttt{client/dist/assets/books/}、クイズの 正解は サーバーだけ。
\end{itemize}

\subsection{F3　こども用ホーム「きょうの がくしゅう」}
\begin{itemize}
\item ログインしたら、先生が 出した 課題・つぎに 読む 本・ロードマップの つぎの 一歩を \textbf{大きな カード 1〜3 まい}で 出す（中身は 宿題の ピン・タスク・ロードマップ。新しい 判定は つくらない）。
\item カードを 押すと その 場所へ \textbf{自動で 移動して ひらく}（F8）。
\item メイン画面の ボタンは こども用では \textbf{5 個 前後}（きょうの がくしゅう・本だな・しまの ちず・きせかえ・ホーム）に しぼり、ほかは「もっと 見る」へ。先生・上級者は いまの まま。
\item 「課題を おえてから 遊びを ひらく」は \textbf{クラスごとの 設定}（先生が オン・オフ。はじめは オフ）。
\item こども用ホームと ボタンの しぼりこみは \textbf{幼児モード（F11）の 中身}として 出す（ふつうの モードでは いまの 画面）。
\end{itemize}

\subsection{F4　全テキストの 読み上げ}
\begin{itemize}
\item こども用ホーム・よみもの島・クイズ画面・宿題画面の \textbf{問題文・選択肢・ボタン名・説明}に 読み上げボタン。押すと 自然な 声で 読む（日本語は 日本語の 声、英語は 英語の 声）。
\item 漢字は やめ、ひらがな＋英語＋アイコンに そろえる（とくに 冒険ノート・復習・交流ラウンジの 見出し）。
\end{itemize}

\subsection{F5　1 画面 1 問の クイズ画面（共通）}
\begin{itemize}
\item 上に 問題番号の 丸（1〜5・いま どこ）、まんなかに 問題文、大きな 選択肢、\textbf{ぜんぶに 読み上げボタン}、右下に「つぎへ／おわり」。正解・不正解の 音は いまの 効果音。
\item まず よみもの島と 宿題に 使い、英単語ハウス・英検の島へ ひろげる（7 形式の 部品の 見た目を そろえる）。
\end{itemize}

\subsection{F6　レベルの 進み具合バー}
\begin{itemize}
\item 「\textbf{いま 段階 A → つぎ 段階 B・あと 3 さつ／12 問}」の バーを こども用ホーム・ロードマップ・\textbf{保護者レポートの 冒頭}に 出す。英検 5 級・4 級の 目標とも 対応させる（ロードマップの 計算を そのまま 使う）。
\end{itemize}

\subsection{F7　クラス参加の 入力を なくす}
\begin{itemize}
\item 先生が 名簿を 登録 → 教室の \textbf{QR コード／URL}（クラス入り・すでに \texttt{?class=} が ある）で 開くと、\textbf{名簿の 名前（アイコン）を えらぶだけ}で 入れる。クラスコードの 手入力を なくす。
\item 先生の 行は いまの とおり 先生パスワードが 要る。名前の 一覧は その クラスの QR から 開いた ときだけ 見せる（よその 子には 見えない）。
\item 手入力の 入り方は 予備として のこす（家の パソコンなど）。
\end{itemize}

\subsection{F8　タップで 移動}
\begin{itemize}
\item 行き先ボタン（「本だなへ」「つりばへ」など）を 押すと \textbf{自動で 移動}（授業モードで 使っている 移動の しくみを 使う）。WASD の 説明は タブレットでは 出さない（いまも 出していない）。
\end{itemize}

\subsection{F9　保護者レポートの 冒頭に「数字 5 つ」}
\begin{itemize}
\item \textbf{学習時間・話した 回数・読んだ 本の 冊数・クイズの 正答率・覚えた 単語}を いちばん 上に。その 下に 進み具合バー（F6）と 最近 読んだ 本（表紙）。web・PDF・LINE の 週／月 レポートに 同じ 形で 出す。
\end{itemize}

\subsection{F10　先生の 課題配信を 3 ステップに}
\begin{itemize}
\item 「\textbf{えらぶ → 生徒を えらぶ → おくる}」。宿題に \textbf{「えほん」の 形}を 1 つ 足し、本だなから 本を えらべる ように する。いまの 宿題は クラスごと なので、\textbf{生徒ごと・グループごとの 配信}を 足す。しめきり・すすみぐあい・保護者レポートの 宿題の 行は いまの まま 使う。
\end{itemize}

\subsection{F11　幼児モード（先生が 管理画面で オン・オフ）}
Raz-Kids なみに かんたんな 画面に きりかえる モード。幼稚園・年長〜小 1・はじめての 子の 教室むけ。\textbf{F3〜F8 を まとめて 1 つの スイッチ}に する。
\begin{longtable}{@{}L{34mm}L{136mm}@{}}
\toprule
\textbf{どこで きりかえる} & 先生の 画面（\texttt{/teacher}・教室の パスワード）と 運営の /admin の 名簿。\textbf{クラスごと}の 既定（オン／オフ）＋ \textbf{子どもごと}に 上書き（名簿の 列「幼児モード」・CSV でも）。かえると 入ったままの 子にも 1 分以内に とどく（名簿の 読みなおしと 同じ しくみ）。\\
\textbf{画面} & こども用ホーム（F3）が はじめに 出る。ボタンは 5 個 前後（きょうの がくしゅう・本だな・しまの ちず・きせかえ・ホーム）。上の バーは コインと 読み上げボタン だけ。冒険ノート・ランキング・交流ラウンジ・チャットの 自由入力・ビデオ通話・同梱ゲームは 出さない（先生が 個別に ゆるせる）。\\
\textbf{文字と 声} & 説明は ひらがなと 絵だけ。ぜんぶ 読み上げボタン つき（F4）。英語は 大きく 1 行。説明の 文は 1 画面 15 字 前後まで。\\
\textbf{操作} & タップだけ（F8 の 行き先ボタンで 自動移動）。スティック・WASD・一人称は 出さない。カメラは 自動で ついてくる。\\
\textbf{学習} & クイズは 1 画面 1 問（F5）・2〜3 択・正解まで やりなおし。本だなは ロードマップの A〜C と 1 の 段階から。英単語ハウスは SUPER EASY。\\
\textbf{島} & 行ける 島を しぼる（メインの島・よみもの島・つり・きせかえ・じぶんの 島）。ほかの 島は 先生が ゆるした ものだけ。\\
\textbf{安全} & 自由入力の チャットは 出さない（定型文だけ）。ほかの 子の 名前の 表示は 先生が きめる。\\
\textbf{つかう もの} & 名簿の 列（プレミアム・停止と 同じ 形）・名簿の 読みなおし（\texttt{recheckRole}）・画面の 出し分けは \texttt{body} の しるし（\texttt{data-mode=kids}）と CSS（スマホ・PC の 出し分けと 同じ やり方）。判定・コイン・記録は ふつうの モードと 同じ（数字が くらべられる）。\\
\bottomrule
\end{longtable}

\section{既存機能との 対応表（新しく つくるのは 太字だけ）}
\begin{longtable}{@{}L{48mm}L{122mm}@{}}
\toprule
\textbf{要件} & \textbf{使う もの（すでに ある）／新しく つくる もの} \\ \midrule
島・建物・地図・光の柱 & island-kit・HUBS／ACTIVITY\_HUBS・しまの ちずの「おすすめ」\\
読み上げ・ゆっくり・光る ことば & /api/tts（自然な 声・キャッシュ）・speech.js・ウーピー。\textbf{ことばを 光らせる 時間わり}\\
音読の 判定 & 面接・話すの館の 音声認識・judgeSpoken・読めた 語が 光る 音読画面 \\
クイズ & 問題DB（qbank）・formats-core／formats-ui・効果音。\textbf{1 画面 1 問の 共通枠} \\
段階・上がる 条件・英検の 見込み & ロードマップ（roadmap.js）・学習記録（roblox\_events） \\
ごほうび & U-Speak コイン・きせかえ・じぶんの 島・ぶきや・ロードマップの スタンプ・ランキング（指標を 足す）\\
ことば & たんごちょう・苦手単語の 復習 \\
課題 & 宿題（/teacher/homework）。\textbf{「えほん」の 形・生徒ごとの 配信} \\
先生 & /class（ロードマップ・表・メモ）・授業モード（Move together・Review）\\
保護者 & web レポート・LINE 週／月。\textbf{冒頭の 数字 5 つ・よんだ 本} \\
クラス参加 & 名簿・?class= の URL・先生パスワード。\textbf{名前を えらぶ 画面・QR} \\
こども用ホーム & 宿題の ピン・タスク・ロードマップ。\textbf{きょうの がくしゅう 画面・ボタンの しぼりこみ} \\
幼児モード & 名簿の 列・名簿の 読みなおし・CSS の 出し分け。\textbf{先生画面の スイッチ・\texttt{data-mode=kids} の 画面} \\
絵本 & \textbf{絵本の 画面（きく・よむ）・絵本の 材料づくり・本の 一覧の 管理ページ} \\
\bottomrule
\end{longtable}

\section{データと サーバー}
\begin{itemize}
\item 本：\texttt{client/dist/assets/books/<id>/}（本文・絵・音声）。クイズ：問題DB（形式「えほん」・級＝段階の 級）。
\item 子どもの 読書記録：\texttt{DATA\_DIR/reading.json}（本ごとの ✓・★・正確さ・語数・時間）。学習記録には world=reading で 書く。
\item 録音：\texttt{DATA\_DIR/recordings/}。§6 の きまりで 保存・削除。
\item 新しい 部屋メッセージ：\texttt{read:shelf}・\texttt{read:open}・\texttt{read:step}・\texttt{read:quiz}・\texttt{read:answer}・\texttt{read:aloud}（判定・コインは 部屋）。
\end{itemize}

\section{非機能要件}
\begin{itemize}
\item \textbf{端末}：iPad（Safari）・Android タブレット・スマホ横もち・PC。古い iPad でも 絵本の 画面は 2D で 軽く（WebGL を 使わない）。絵は 1 ページずつ 読みこみ（つぎの ページだけ 先読み）、1 枚 200KB 以下。
\item \textbf{通信}：弱い Wi-Fi でも 読めるよう 音声は 先に 作って キャッシュ（同じ 文は 2 回目から ファイル）。
\item \textbf{録音}：先生だけ 聞ける・保存 30 日（設定で かえる）・1 人ぶんの 削除に ふくめる・録音しない 設定（文字の 判定だけ）も もつ。
\item \textbf{費用}：絵と 音声は 1 回 作って 保存（何度 読んでも 追加費用なし）。AI の 物語・絵の 生成は 本ごとに 1 回。
\end{itemize}

\section{開発の 段階と 受け入れ基準}
\begin{longtable}{@{}L{18mm}L{72mm}L{80mm}@{}}
\toprule
\textbf{段階} & \textbf{中身} & \textbf{受け入れ基準（できたと 言える 条件）} \\ \midrule
すぐ & F4 全テキスト読み上げ・F7 入力なしの クラス参加・F8 タップで 移動 & 初めての 子が \textbf{文字を 1 字も 打たずに} QR から 入れる。こども用の 画面の 文字は ぜんぶ 読み上げボタンで 聞ける。行き先ボタンで 目的地に 着く。\\
つぎ & F3 こども用ホーム・F11 幼児モード（先生の スイッチ）・F5 1 画面 1 問・F6 進み具合バー・F9 保護者の 数字 5 つ & 先生が 管理画面で クラス／子どもごとに 幼児モードを オンに でき、1 分以内に 画面が かわる。幼児モードでは ログイン後 1 タップで 今日の 学習が はじまり、ボタンは 5 個 前後・文字を 打つ 場面が ない。保護者レポートと LINE の いちばん 上に 数字 5 つと 進み具合バー。\\
本命 & F1 よみもの島・F2 絵本 36 冊（→120 冊）・F10 3 ステップ配信 & 12 段階 × 3 冊が 読める。1 冊を きく→よむ→おんどく→クイズ で ✓・コイン・学習記録に のる。先生が 3 ステップで 本を 生徒ごとに 出せる。\\
\bottomrule
\end{longtable}
\textbf{検査}：サーバーの テスト（判定・コイン 1 回・記録・権限）、ブラウザの テスト（iPad／スマホ横もちの 画面・11px・はみ出し・英語モードで 学習の 中身が 訳されない）、実機の iPad で 読み上げと マイクを 確認。

\section{きめて ほしい こと（承認の とき いっしょに）}
\begin{enumerate}
\item 段階の 順番：「すぐ → つぎ → 本命」で よいか（よみもの島を 先に したい 場合は 本命から）。
\item 島の 名前（案：\textbf{よみもの島／Book Island}）。
\item 絵本の 材料：AI オリジナル中心で よいか・大学の 素材は いつ 何冊 とどくか。
\item 音読の 録音を のこすか（のこすなら 30 日で よいか）。
\item 「課題を おえてから 遊びを ひらく」を クラス設定に するか（はじめは オフ）。
\item 幼児モードで 出す 島と ボタンの 一覧（案は F11）・既定を オンに する クラス（幼稚園の 教室など）。
\item クラス参加で 名簿の 名前を 一覧で えらばせて よいか（QR から 開いた ときだけ）。
\end{enumerate}

\begin{note}
\textbf{この 要件定義に「Yes」を いただくまで、実装は はじめません。} 変えたい ところが あれば この 文書を 直して 版を 上げます。\\
出典：「Raz-Kids 徹底リサーチ」（U-Speak Lab, 2026-11-17）、Learning A-Z 公開資料（Raz-Kids 概要・Assessing a Student's Level）、Common Sense Education、U-Speak 商談記録。
\end{note}

\end{document}
